\begin{minipage}[t]{0.5\linewidth}
Le but de l'exercice est d'étudier un exemple de droite d'Euler.

On considère un cube ABCDEFGH de côté une unité.

L'espace est muni du repère orthonormé $\left(\mathrm{A} ~;~ \vectt{AB} ~;~ \vectt{AD} ~;~ \vectt{AE}\right)$.

On note I le milieu du segment [AB] et J le milieu du segment [BG].
\end{minipage}
\hfill
\begin{minipage}[t]{0.45\linewidth}
\psset{unit=0.7cm, radius=0pt}
\def\xmin{-2} \def\xmax{6.5} \def\ymin{-0.5} \def\ymax{5.5}
\begin{pspicture}(\xmin,\ymin)(\xmax,\ymax)

\Cnode*(0,0){B} \Cnode*(4,0){C} \Cnode*(0,4){F} \Cnode*(4,4){G}
\Cnode*(2,1){A} \Cnode*(6,1){D} \Cnode*(2,5){E} \Cnode*(6,5){H}
\psline(F)(G)(C)(B)(F)(E)(H)(D)(C)
\psline(G)(H)
\psline[linestyle=dashed](D)(A)(B) \psline[linestyle=dashed](E)(A)

\small
\uput[ul](A){A} \uput[dl](B){B} \uput[dr](C){C} \uput[r](D){D}
\uput[u](E){E} \uput[ul](F){F} \uput[dr](G){G} \uput[ur](H){H}


\end{pspicture}
\end{minipage}

\begin{enumerate}
\item On donne les coordonnées des points A, B, G, I et J:

A\;$\left (0 \;,\; 0 \;,\;0 \right )$; B\;$\left (1 \;,\; 0 \;,\; 0 \right )$; G\;$\left ( 1\;,\;  1\;,\;1  \right )$; I\;$\left ( \dfrac{1}{2} \;,\; 0 \;,\; 0 \right )$; J\;$\left ( 1 \;,\; \dfrac{1}{2} \;,\;  \dfrac{1}{2}\right )$
\item
	\begin{enumerate}
		\item 
La droite (AJ) est l'ensemble des points M\,$\left (x\;,\; y \;,\; z\right )$ tels que $\vectt{AM}$ et $\vectt{AJ}$ soient colinéaires, c'est-à-dire tels que $\vectt{AM} = k.\vectt{AJ}$ avec $k\in\R$.

$\vectt{AM}$ a pour coordonnées  
$\begin{pmatrix}
x_{\text{M}} - x_{\text{A}} \\ y_{\text{M}} - y_{\text{A}} \\ z_{\text{M}} - z_{\text{A}}
\end{pmatrix}
=
\begin{pmatrix}
x\\ y \\ z
\end{pmatrix}$
et celles de $\vectt{AJ}$ sont
$\begin{pmatrix}
x_{\text{J}} - x_{\text{A}} \\ y_{\text{J}} - y_{\text{A}} \\ z_{\text{J}} - z_{\text{A}}
\end{pmatrix}
=
\begin{pmatrix}
1\\ \dfrac{1}{2} \\[5pt] \dfrac{1}{2}
\end{pmatrix}$

$\vectt{AM} = k.\vectt{AJ}
\iff
\begin{pmatrix}
x\\ y \\ z
\end{pmatrix}
=
k.\begin{pmatrix}
1\\ \dfrac{1}{2} \\[5pt] \dfrac{1}{2}
\end{pmatrix}
\iff
\left \{
\begin{array}{l !{=} l}
x & k\\
y & \dfrac{k}{2}\\[5pt]
z & \dfrac{k}{2}
\end{array}
\right .$

La droite (AJ) a donc pour représentation paramétrique
$\left \{
\begin{array}{l !{=} l }
x & k \\
y & \dfrac{k}{2} \\[5pt]
z & \dfrac{k}{2}
\end{array}
\right. \text{avec } k\in \R
$
\item 
$\left\{
\begin{array}{l !{=} l }
x & \dfrac{1}{2}+\dfrac{1}{2} t \\[5pt]
y &  t \\
z & t\end{array}
\right. (t\in\R)$
est la représentation paramétrique d'une droite passant par le point de coordonnées 
$\left( \dfrac{1}{2}\;,\; 0 \;,\; 0\right )$, c'est-à-dire le point I, et  ayant le vecteur de coordonnées 
$\begin{pmatrix}
\dfrac{1}{2} \\[5pt] 1 \\ 1
\end{pmatrix}$
pour vecteur directeur.

$\vectt{IG}$ a pour coordonnées 
$\begin{pmatrix}
x_{\text{G}} - x_{\text{I}} \\ y_{\text{G}} - y_{\text{I}} \\ z_{\text{G}} - z_{\text{I}}
\end{pmatrix}
=
\begin{pmatrix}
1-\dfrac{1}{2}\\[5pt] 1-0 \\ 1-0
\end{pmatrix}
=
\begin{pmatrix}
\dfrac{1}{2}\\[5pt] 1 \\ 1
\end{pmatrix}$

Donc $\left\{
\begin{array}{l !{=} l }
x & \dfrac{1}{2}+\dfrac{1}{2} t  \\[5pt]
y & t \\
z & t
\end{array} 
\right.\; (t\in\R)$
est une représentation paramétrique de la droite (IG).

			\item
Les droites (AJ) et (IG) sont sécantes si et seulement si on peut trouver un couple de réels $\left (k\;,\;t\right )$ tel que
$\left\{
\begin{array}{l !{=} l }
k & \dfrac{1}{2}+\dfrac{1}{2} t \\[5pt]
\dfrac{1}{2}k & t \\[5pt]
\dfrac{1}{2}k & t
\end{array} 
\right.$.
On résout ce système.

$\begin{aligned}[t]
\left\{
\begin{array}{l !{=} l }
k & \dfrac{1}{2}+\dfrac{1}{2} t \\[5pt]
\dfrac{1}{2}k & t \\[5pt]
\dfrac{1}{2}k & t
\end{array}
\right.
&
\iff
\left\{
\begin{array}{l !{=} l }
k & \dfrac{1}{2}+\dfrac{1}{2} t \\[5pt]
k & 2 t
\end{array} 
\right.
\iff
\left\{
\begin{array}{l !{=} l }
2t & \dfrac{1}{2}+\dfrac{1}{2} t \\[5pt]
k & 2 t
\end{array} 
\right.
\iff
\left\{
\begin{array}{l !{=} l }
\dfrac{3}{2} t & \dfrac{1}{2} \\[5pt]
k & 2 t
\end{array}
\right.\\
&
\iff
\left\{
\begin{array}{l !{=} l }
 t & \dfrac{1}{3} \\[5pt]
k & \dfrac{2}{3}
\end{array}
\right.
\end{aligned}$

Donc les droites (AJ) et (IG) sont sécantes.
Pour $k=\dfrac{2}{3}$, on a
$\left \{
\begin{array}{l !{=} l }
k & \dfrac{2}{3} \\[5pt]
\dfrac{1}{2}k &  \dfrac{1}{3} \\[5pt]
\dfrac{1}{2}k &  \dfrac{1}{3}
\end{array} 
\right.$

Le point S d'intersection des droites (AJ) et (IG) a donc pour coordonnées
$\left (\dfrac{2}{3}\;,\; \dfrac{1}{3} \;,\; \dfrac{1}{3}\right )$.
		\end{enumerate}

		\item \begin{enumerate}
			\item 
Soit $\vect{n}$ le vecteur de coordonnées $(0 ~;~-1 ~;~ 1)$.

$\vectt{AB}$ a pour coordonnées 
$\begin{pmatrix}
1\\ 0 \\ 0
\end{pmatrix}$
et $\vectt{AG}$ a pour coordonnées
$\begin{pmatrix}
1\\ 1\\1
\end{pmatrix}$

\begin{list}{\textbullet}{}
\item $\vect{n}\cdot \vectt{AB} = 0\times 1 + (-1)\times 0 +1\times 0=0$ donc $\vect{n} \perp \vectt{AB}$.
\item $\vect{n}\cdot \vectt{AG} = 0\times 1 + (-1)\times 1 +1\times 1=0$ donc $\vect{n} \perp \vectt{AG}$.
\item Les vecteurs $\vectt{AB}$ et $\vectt{AG}$ ne sont pas colinéaires donc ce sont deux vecteurs directeurs du plan (ABG).
\end{list}

Le vecteur $\vect{n}$ est orthogonal à deux vecteurs directeurs du plan (ABG), donc $\vect{n}$ est un vecteur normal au plan (ABG).

		\item 
Le plan (ABG) est l'ensemble des points M de coordonnées $(x\;,\; y \;,\; z)$ tels que $\vectt{AM}$ et $\vect{n}$ soient orthogonaux, c'est-à-dire tels que $\vectt{AM}\cdot \vect{n}=0$.

$\vectt{AM}\cdot \vect{n}=0
\iff
x \times 0+ y\times (-1) +z\times 1 = 0
\iff
-y+z=0$

Le plan (ABG) a pour équation cartésienne: $-y+z=0$.

		\item On admet qu'une représentation paramétrique de la droite $(d)$ de vecteur directeur $\vect{n}$ et passant par le point K de coordonnées $\left(\dfrac{1}{2} ~;~ 0 ~;~ 1\right)$ est : $\left\{\begin{array}{l}
x=\dfrac{1}{2} \\[5pt]
 y= - t \\
  z=1+t \end{array} \right.\; (t \in \R).$



S'il existe, le point d'intersection de la droite $(d)$ et du plan (ABG), est solution du système:
$\left \{
\begin{array}{r !{=} l}
x & \dfrac{1}{2} \\[5pt]
y & -t\\
z & 1+t\\
-y+z & 0
\end{array}
\right .$

Il faut donc que: $-(-t)+(1+t)=0$ soit $2t+1=0$ ou encore $t=-\dfrac{1}{2}$.

Pour $t=-\dfrac{1}{2}$, on a $x=\dfrac{1}{2}$, $y=-t = \dfrac{1}{2}$ et $z= 1-\dfrac{1}{2}=\dfrac{1}{2}$.

Donc la droite $(d)$ coupe le plan (ABG) au point L de coordonnées $\left ( \dfrac{1}{2}\;,\; \dfrac{1}{2}\;,\; \dfrac{1}{2}\right )$.

			\item 
\begin{list}{\textbullet}{On calcule les distances LA, LB et LG.}
\item
$\begin{aligned}[t]
\text{LA}^2  = \left (x_{\text{A}} - x_{\text{L}}\right )^2 
+ \left (y_{\text{A}} - y_{\text{L}}\right )^2 
+ \left (z_{\text{A}} - z_{\text{L}}\right )^2\\
= \left (0 - \dfrac{1}{2}\right )^2 
+ \left (0 - \dfrac{1}{2}\right )^2 
+ \left (0 - \dfrac{1}{2}\right )^2  
= \dfrac{1}{4} + \dfrac{1}{4} + \dfrac{1}{4} 
= \dfrac{3}{4}
\end{aligned}$

donc $\text{LA} = \sqrt{\dfrac{3}{4}} = \dfrac{\sqrt{3}}{2}$.

\item 
$\begin{aligned}[t]
\text{LB}^2 &= \left (x_{\text{B}} - x_{\text{L}}\right )^2 
+ \left (y_{\text{B}} - y_{\text{L}}\right )^2 
+ \left (z_{\text{B}} - z_{\text{L}}\right )^2\\
&= \left (1 - \dfrac{1}{2}\right )^2 
+ \left (0 - \dfrac{1}{2}\right )^2 
+ \left (0 - \dfrac{1}{2}\right )^2 \\
&= \dfrac{1}{4} + \dfrac{1}{4} + \dfrac{1}{4} \\
&= \dfrac{3}{4}
\end{aligned}$

donc $\text{LB} = \sqrt{\dfrac{3}{4}} = \dfrac{\sqrt{3}}{2}$.

\item
$\begin{aligned}[t]
\text{LG}^2 &= \left (x_{\text{G}} - x_{\text{L}}\right )^2 
+ \left (y_{\text{G}} - y_{\text{L}}\right )^2 
+ \left (z_{\text{G}} - z_{\text{L}}\right )^2\\
&= \left (1 - \dfrac{1}{2}\right )^2 
+ \left (1 - \dfrac{1}{2}\right )^2 
+ \left (1 - \dfrac{1}{2}\right )^2 \\
&= \dfrac{1}{4} + \dfrac{1}{4} + \dfrac{1}{4} \\
&= \dfrac{3}{4}
\end{aligned}$
donc $\text{LG} = \sqrt{\dfrac{3}{4}} = \dfrac{\sqrt{3}}{2}$.
\end{list}

Donc le point L est équidistant des points A, B et G.
		\end{enumerate}

		\item 
$\vectt{AB}$ a pour coordonnées  
$\begin{pmatrix}
1\\ 0 \\ 0
\end{pmatrix}$
et $\vectt{BG}$ a pour coordonnées
$\begin{pmatrix}
1 - 1 \\ 1-0 \\1-0
\end{pmatrix}
=
\begin{pmatrix}
0 \\ 1 \\1
\end{pmatrix}$

$\vectt{AB}\cdot \vectt{BG} = 1\times 0 + 0\times 1 + 0\times 1 = 0$
donc $\vectt{AB}\perp \vectt{BG}$ donc le triangle ABG est rectangle en B.

	\item
		\begin{enumerate}
			\item
\begin{list}{\textbullet}{Points remarquables du triangle ABG.}
\item (AJ) et (IG) sont deux médianes du triangle ABG et elles se coupent en S; donc le point S est le centre de gravité du triangle ABG.
\item Le point L est équidistant des points A, B et G donc c'est le centre du cercle circonscrit au triangle ABG.
\item Le triangle ABG est rectangle en B donc B est son orthocentre.
\end{list}
			\item 
$\vectt{BS}$ a pour coordonnées  
$\begin{pmatrix}
\dfrac{2}{3} - 1 \\[5pt] 
\dfrac{1}{3} - 0 \\[5pt] 
\dfrac{1}{3} - 0
\end{pmatrix}
=
\begin{pmatrix}
-\dfrac{1}{3} \\[5pt] 
\dfrac{1}{3}  \\[5pt] 
\dfrac{1}{3} 
\end{pmatrix}$
et  $\vectt{BL}$ a pour coordonnées
$\begin{pmatrix}
\dfrac{1}{2} - 1 \\[5pt] 
\dfrac{1}{2} - 0 \\[5pt] 
\dfrac{1}{2} - 0
\end{pmatrix}
=
\begin{pmatrix}
-\dfrac{1}{2} \\[5pt]
\dfrac{1}{2} \\[5pt]
\dfrac{1}{2} 
\end{pmatrix}$

$\vectt{BS} = \dfrac{2}{3} \vectt{BL}$ donc les vecteurs $\vectt{BS}$ et $\vectt{BL}$ sont colinéaires et donc les points B, S et L sont alignés.

		\end{enumerate}
	\end{enumerate}
